\begin{chaptersummary}
  \item The router is a downloaded binary, not a package, and its default is to
        fetch its execution config from WunderGraph's CDN with a graph token.
        Given none it opens its listeners, announces a playground and a metrics
        endpoint nobody asked for, and then refuses to start. The refusal names
        the local alternative, which is one key in one file. The port this book
        uses, 3002, is the one the router prints on the way down.
  \item Handed a static execution config the router stops polling anything, and
        reads the file once: changing the graph means restarting the process. It
        starts happily with no subgraph running, because loading a config and
        calling a service are different things, the same way composition reads
        files and never calls a service. The nil uuid
        chapter~\ref{ch:composition} found in \code{router.json} comes back out
        here as \code{config\_version}.
  \item What the router serves is the subgraph's schema minus exactly the
        federation route: three types and two \code{Query} fields, including
        \code{\_entities}. Hidden is not closed. The subgraph still answers
        \code{\_service} and \code{\_entities} on its own port, so the router's
        schema describes one of your two surfaces.
  \item A query plan is an Advanced Request Tracing option, so without a Cosmo
        Cloud token the header asking for one is ignored in silence rather than
        refused. \code{dev\_mode: true} turns it on locally and the router then
        explains that it has. Every plan over this graph is a \code{Sequence}
        holding one \code{Single} fetch, and the router costs the graph nothing:
        two statements through it, two statements without it, none at all for a
        plan asked for with \code{X-WG-Skip-Loader}. Reading a plan also gives
        away that the router rewrites literal arguments into variables before
        forwarding, and that it answers introspection itself rather than asking
        anyone.
  \item A subgraph that is not running produces HTTP 200, a null at the first
        nullable field, and an error naming which service failed and nothing
        about why. Both health endpoints go on answering OK throughout, which is
        correct, because they report that the router is up rather than that the
        graph works. Those are different claims and only one of them is being
        made.
\end{chaptersummary}
